\documentclass[12pt,a4paper]{article}
\usepackage{fontawesome5}
\providecommand{\faPencilAlt}{\faIcon{pencil-alt}}
\usepackage{styles/couleurs}
\usepackage{styles/mise_en_page}
\usepackage{styles/boites}
\usepackage{styles/maths}
\usepackage{styles/figures}
\usepackage{enumitem}
\usepackage{array}
\usepackage{tabularx}
\usepackage{amsmath}
\usepackage{amssymb}
\usepackage{tikz}
\usepackage{xparse}
\usetikzlibrary{calc,patterns,arrows.meta}
\setlist[itemize]{leftmargin=1.4em,itemsep=0.15em,topsep=0.2em}
\setlist[enumerate]{leftmargin=*,itemsep=0.25em,topsep=0.25em}
\newcommand{\lignepointillee}{\noindent\dotfill\par\vspace{0.45em}}
\newcommand{\lignesreponse}[1]{%
\begin{tcolorbox}[enhanced,colback=white,colframe=NeutreC!40,boxrule=0.5pt,
arc=1.5mm,left=3mm,right=3mm,top=2mm,bottom=1mm,before skip=2pt,after skip=2pt]
\foreach \i in {1,...,#1}{\lignepointillee}
\end{tcolorbox}}
\newcommand{\titreseq}[1]{%
\begin{center}
{\sffamily\LARGE\bfseries\color{TitreC}Séquence 4 -- Calcul littéral}\\[0.2em]
{\sffamily\large\bfseries\color{SavoirC}#1}
\end{center}}
\begin{document}

\titreseq{Activité 7 -- Une conjecture est-elle une preuve ?}

\section*{Partie A -- Un programme qui semble toujours revenir au départ}

\begin{tcolorbox}[colback=white,colframe=SavoirC,boxrule=0.7pt,arc=2mm]
\begin{enumerate}
  \item Choisir un nombre.
  \item Le doubler.
  \item Ajouter \(5\).
  \item Doubler le résultat.
  \item Retirer le triple du nombre de départ.
  \item Retirer \(10\).
\end{enumerate}
\end{tcolorbox}

\begin{enumerate}
  \item Applique le programme à \(2\), \(-3\) et \(10\).
  \item Énonce une conjecture.
  \item Trois essais constituent-ils une preuve ?
\end{enumerate}

\lignesreponse{7}

\section*{Partie B -- Prouver avec une lettre}

On note \(x\) le nombre de départ.

\begin{enumerate}
  \item Écris l'expression traduisant le programme.
  \item Développe puis réduis.
  \item Explique pourquoi le résultat constitue une preuve valable
  pour n'importe quel nombre de départ.
\end{enumerate}

\lignesreponse{8}

\section*{Partie C -- Le programme trompeur}

\begin{tcolorbox}[colback=white,colframe=ActionC,boxrule=0.7pt,arc=2mm]
\begin{enumerate}
  \item Choisir un nombre.
  \item Ajouter \(1\) à son carré.
  \item Multiplier par \(6\).
  \item Retirer le cube du nombre de départ.
  \item Diviser par \(11\).
\end{enumerate}
\end{tcolorbox}

Applique ce programme à \(1\), \(2\), \(3\), puis \(4\).

\begin{enumerate}
  \item Que pourrait-on conjecturer après les trois premiers essais ?
  \item Que montre l'essai avec \(4\) ?
  \item Complète : un exemple peut \dotfill{} une affirmation générale,
  mais plusieurs exemples ne suffisent pas à la \dotfill
\end{enumerate}

\lignesreponse{7}

\aretenir{Pour prouver que deux expressions sont égales pour toute valeur de la lettre,
on peut les transformer jusqu'à obtenir la même expression.

Pour montrer qu'une affirmation générale est fausse, un seul contre-exemple suffit.}

\end{document}
